%% kapa: exam=1 round=19 year=2008 period=2 form=A subject=english subject_ko=영어 seq=24 number=64
%% source: Q-Net 감정평가사 기출문제 (cst00309) artlSeq=1001962 file=extracted/1008/1001962/2008년 제19회 감정평가사 1차시험 기출문제/2교시.hwp
%% conversion: kordoc 4.24.51 -> tools/gen1.py (auto); review=proofread-sonnet; flags=none
\begin{question}{64}
% [공통 지문/지시문: 원문에서 여러 문항에 걸쳐 제시됨]
\begin{framedtext}
※ 밑줄 친 (A), (B), (C)에 들어갈 단어로 가장 알맞게 짝지어진 것을 고르시오. (63-64) \par
\end{framedtext}
Learning (A) \uline{what / how} to budget and save money should start early. Living within our means is a lesson that should be (B) \uline{learned / learning} at a young age and  reinforced throughout our lives.  Staying away from debt is better than (C) \uline{get / getting} out of debt. But in our consumer society, it is not easy.
\begin{choices}
  \item[] \makebox[5em][l]{(A)}\makebox[1.5em]{}\makebox[5em][l]{(B)}\makebox[1.5em]{}(C)
  \item \makebox[5em][l]{what}\makebox[1.5em][c]{―}\makebox[5em][l]{learned}\makebox[1.5em][c]{―}get
  \item \makebox[5em][l]{what}\makebox[1.5em][c]{―}\makebox[5em][l]{learning}\makebox[1.5em][c]{―}getting
  \item \makebox[5em][l]{how}\makebox[1.5em][c]{―}\makebox[5em][l]{learned}\makebox[1.5em][c]{―}get
  \item \makebox[5em][l]{how}\makebox[1.5em][c]{―}\makebox[5em][l]{learned}\makebox[1.5em][c]{―}getting
  \item \makebox[5em][l]{how}\makebox[1.5em][c]{―}\makebox[5em][l]{learning}\makebox[1.5em][c]{―}get
\end{choices}
\end{question}
\answer{4}
